%==============================================================================
%  IEEE TRANSACTIONS ON INFORMATION FORENSICS AND SECURITY (TIFS)
%  Manuscript skeleton -- DeployBench: on-device benchmark of speech CMs
%
%  Structure follows 00_planning/paper_outline.md (v2). Section-level guidance
%  is left in % comments; each cites the governing decision (Dxx) or doc.
%  Metric defs encode the Tier-1 reconciliation (TTG ref term = min-HTER, D19
%  candidate; QTS per experiment_protocol.md sec.3.5).
%
%  Compile: pdflatex -> bibtex -> pdflatex x2   (needs refs.bib alongside)
%  Template: standard IEEEtran journal class (TIFS is NOT compsoc).
%==============================================================================
\documentclass[journal]{IEEEtran}

% ---- packages -------------------------------------------------------------
\usepackage{amsmath,amssymb,amsfonts}
\usepackage{algorithmic}
\usepackage{graphicx}
\usepackage{booktabs}
\usepackage{multirow}
\usepackage{siunitx}          % latency/energy units
\usepackage{xcolor}
\usepackage[hidelinks]{hyperref}
\usepackage{url}
\usepackage[capitalize,noabbrev]{cleveref}

% ---- draft helpers (delete before submission) -----------------------------
\newcommand{\todo}[1]{\textcolor{red}{[\textbf{TODO:} #1]}}
\newcommand{\placeholder}[1]{\textcolor{gray}{\emph{#1}}}

% ---- project macros -------------------------------------------------------
\newcommand{\bench}{\textsc{DeployBench}}
\newcommand{\ttg}{\mathrm{TTG}}
\newcommand{\qts}{\mathrm{QTS}}
\newcommand{\hter}{\mathrm{HTER}}
\newcommand{\frr}{\mathrm{FRR}}
\newcommand{\far}{\mathrm{FAR}}
\newcommand{\eer}{\mathrm{EER}}
\newcommand{\taustar}{\tau^{\star}}

\begin{document}

% ============================ TITLE ========================================
% Working title (paper_outline.md). Shorten for the running head.
% RETITLE BEFORE SUBMISSION (D30). With the roster frozen at five, "Benchmark"
% in the title invites the breadth objection at first glance -- lm-Meter,
% WhisperKit and MLPerf all cover far more systems. The paper is a PROTOCOL
% paper with a demonstration audit; lead with the protocol and the two metrics.
\title{Can Deepfake Voice Detectors Run on Your Phone?\\
An End-to-End On-Device Benchmark of Speech\\ Anti-Spoofing Countermeasures}

\author{
  Hoan~Hoang,~\IEEEmembership{}
  and~\todo{co-authors, advisor}
  \thanks{Manuscript submitted \todo{Jan 2027}. \todo{funding/affiliations.}}
  \thanks{Code and harness: \url{https://github.com/\todo{repo}}.}
}

% Running heads
\markboth{IEEE Transactions on Information Forensics and Security,~Vol.~XX, No.~X, \todo{Month Year}}%
{Hoang \MakeLowercase{\textit{et al.}}: On-Device Benchmark of Speech Anti-Spoofing Countermeasures}

\maketitle

% ============================ ABSTRACT =====================================
\begin{abstract}
% Write LAST. Compress to: (1) the regime gap -- CMs are scored at oracle
% thresholds, fp32, on server GPUs; deployment is fixed tau, quantized, phone
% CPU, lossy channels. (2) What we do -- convert N published CMs through
% reproduce->convert->quantize->deploy, measure on 3 phones under fixed-tau.
% (3) Two coined metrics (TTG, QTS). (4) The two headline findings, with
% numbers, once known. (5) The released harness + deployability guidance.
% Do NOT claim "nobody runs detection on phones" (D2) -- claim no OPEN,
% reproducible, systematic on-device evaluation exists.
\placeholder{One-paragraph abstract, \(\sim\)200 words. Fill after Results freeze.}
\end{abstract}

\begin{IEEEkeywords}
Audio deepfake detection, speech anti-spoofing, countermeasures, on-device
inference, quantization, threshold calibration, half-total error rate,
benchmarking, mobile.
\end{IEEEkeywords}

% ============================ 1. INTRODUCTION ==============================
\section{Introduction}
\label{sec:intro}
\IEEEPARstart{S}{peech} deepfake countermeasures\todo{open the regime gap}.
% ~1.5 pp. Argument spine (paper_outline sec.1):
%  - Published CMs are evaluated at an ORACLE threshold on the labeled test set,
%    in fp32, on server GPUs. Deployment shares none of it: tau fixed in advance,
%    weights quantized, phone CPU, lossy channels. Cite zhou26 for the
%    operating-point blowup (0.21% EER -> 39.5% HTER, 78.7% bona fide rejected).
%  - The generalization gap is documented (muller22, berisha25); the DEPLOYMENT
%    gap end-to-end is not measured anywhere.
%  - Lead with the two headline findings once known (candidates: H1
%    anti-correlation / H3 QTS / worst conversion casualty).

\subsection*{Contributions}
% paper_outline sec.1. Keep the non-goals explicit (D1: no new model).
\begin{itemize}
  \item \textbf{\bench{}}: a reproducible protocol and released harness for
        fixed-threshold, on-device evaluation of speech CMs --- threshold
        transfer, \(\frr/\far\) decomposition, sensitivity curves, a channel
        stress suite, and mobile latency/energy/thermal methodology.
  % OPEN [U] (OD-4): the device grid is ORT-only until OD-4 is decided (decision.md
  % 2026-08-28); "two runtimes" stands only if an ExecuTorch leg lands.
  \item \textbf{An audit} of five countermeasures across
        \(\{\text{fp32},\text{fp16},\text{int8}\}\) precisions, two runtimes,
        four evaluation sets, and 3 physical phones, with pre-registered
        hypotheses and bootstrap confidence intervals.
  \item \textbf{Two metrics}: the \emph{threshold transfer gap} (\ttg{}) and the
        \emph{quantization-induced threshold shift} (\qts{}), the latter not
        previously reported.
  \item \textbf{Deployability guidance}: a per-device-tier decision table and a
        released Android application.
\end{itemize}
\noindent\emph{Non-goals.} We propose no new detector (D1) and do not study
open-world generalization beyond the transfer targets used here.

% ============================ 2. RELATED WORK ==============================
\section{Related Work}
\label{sec:related}
% ~1.5 pp, 6 paragraphs, each a DIFFERENT differentiator (D7). Positioning
% table after para 3. Cites keyed to refs.bib. Corrections from the Tier-1/2
% reconciliation are baked in -- do not "restore" the removed claims.

\subsection{Detection SOTA and its evaluation regime}
Progress in speech anti-spoofing has been organized largely by the ASVspoof
challenge series \cite{yamagishi21_asvspoof2021,wang24_asvspoof5}, and the
current state of the art is a family of self-supervised (SSL) front-ends ---
wav2vec\,2.0, XLS-R, and WavLM \cite{babu22_xlsr,chen22_wavlm} --- coupled to
graph-attention or attentive-pooling back-ends
\cite{jung22_aasist,tak22_xlsr_aasist}. Such systems report equal-error rates
(EER) below \SI{0.3}{\percent} on the ASVspoof\,2019 LA benchmark and dominate
recent leaderboards \cite{li25_survey}. Two properties matter here. First, the
headline systems are large: the SSL back-bones alone span roughly 90 to 300
million parameters, placing their memory and compute well outside a phone's
per-inference budget. Second, they are compared almost exclusively by EER and
detection-cost functions evaluated at an \emph{oracle} operating point chosen on
the labeled test set. We therefore treat these models not as baselines to beat
but as a regime to characterize: we retain one representative SSL system
(XLS-R\,+\,AASIST) as an unranked server-side reference and ask what becomes of
the rest once they are quantized and fixed to a single threshold. Our
contribution here is one of \emph{regime}, not accuracy.

\subsection{Generalization failure}
Cross-corpus fragility has been documented for a decade. M\"{u}ller et
al.\ \cite{muller22_in_the_wild} introduced the In-the-Wild corpus and showed
that a roster of twelve \emph{pre-SSL} architectures loses one to three orders
of magnitude of accuracy on it; SSL front-ends have since narrowed, though not
closed, this gap. Berisha et al.\ \cite{berisha25_coverage_debt} frame the
residual failure structurally: the space of deployment conditions grows
multiplicatively while data grows additively, so a non-zero error floor is
unavoidable and detectors are best used as auxiliary signals rather than sole
gates. The community's response has been overwhelmingly \emph{training-side} ---
new augmentation schemes, re-synthesized hard negatives, and larger multilingual
corpora \cite{mlaad24}.% \todo{method-wave cluster: 2501.14240 / 2509.00186 / 2606.30791}
These interventions change what a model learns but continue to report a single
oracle EER; ours is orthogonal and lies entirely on the evaluation side, where
the threshold is fixed before the target is seen.

\subsection{Thresholds and calibration: the key precedent}
A deployed detector enjoys no oracle: its threshold must be fixed in advance and
applied to unlabeled traffic. Biometric authentication long distinguished
a-priori from a-posteriori thresholds and reported the half-total error rate
(HTER) at a fixed operating point \cite{bengio04_hter}; the tandem detection cost
function \cite{kinnunen18_tdcf} and its ASVspoof\,5 successors refined
cost-sensitive scoring, yet EER remains the de facto headline. The precedent
closest to ours is Zhou and Wang \cite{zhou26_eer_hides}, who freeze a
state-of-the-art CM and transfer its source threshold: a \SI{0.21}{\percent}-EER
detector rejects \SI{78.7}{\percent} of genuine In-the-Wild speech, and the
standard unlabeled corrections either provably cannot alter the ranking or fail
with opposite sign on two targets that differ only in class prior. ASVspoof\,5
independently surfaced calibration as a deployment concern --- its strongest open
submission attains near-oracle minDCF yet a degenerate actDCF
\cite{wang24_asvspoof5}. Benchmark tooling such as AUDDT \cite{auddt25}
aggregates detectors but remains desktop-bound and oracle-metric. Where this line
\emph{diagnoses} that post-hoc score repair is limited, we take the diagnosis as
a design constraint: we revive fixed-threshold HTER for countermeasures,
mandate its \(\frr/\far\) decomposition, and carry it onto physical devices.

% Ported 2026-09-28 from 01_references/related_work_matrix.md (post-reconciliation
% cells; see notes/tier1/_RECONCILIATION_tier1.md). Rows trimmed to works cited
% in this section. One deliberate change from the matrix: the SSL row's channel
% cell is "--", not "partial", because training-time augmentation (RawBoost) is
% not a channel-stress evaluation -- the same rule the Pascu note applies to
% pascu25. Berisha (recommendation only, every cell absent) and AUDDT (channel
% coverage unverified) are cited in text but have no row.
\begin{table}[!t]
\caption{Positioning against prior work. \checkmark: addressed;
$\circ$: partial; --: absent; n/a: not applicable to the task.
\emph{Channel} means evaluation under controlled codec or transmission
degradation, not training-time augmentation. \emph{Harness} means a released,
end-to-end evaluation pipeline; releasing code or data alone is partial.}
\label{tab:positioning}
\centering
\resizebox{\columnwidth}{!}{%
\begin{tabular}{lccccc}
\toprule
Work & Fixed-$\tau$ & On-device & Quant. & Channel & Harness \\
\midrule
AASIST, XLS-R\,+\,AASIST \cite{jung22_aasist,tak22_xlsr_aasist}
  & -- & -- & -- & -- & $\circ$ \\
ASVspoof\,2021 DF \cite{yamagishi21_asvspoof2021}
  & -- & -- & -- & \checkmark & $\circ$ \\
ASVspoof\,5 \cite{wang24_asvspoof5}
  & $\circ$\textsuperscript{a} & -- & -- & $\circ$ & $\circ$ \\
In-the-Wild \cite{muller22_in_the_wild}
  & -- & -- & -- & -- & $\circ$ \\
Zhou and Wang \cite{zhou26_eer_hides}
  & \checkmark & -- & -- & $\circ$\textsuperscript{b} & $\circ$ \\
Pascu et al.\ \cite{pascu25_edge_plugin}
  & -- & $\circ$\textsuperscript{c} & -- & -- & $\circ$ \\
Hybrid pruning \cite{pruning_ssl}
  & -- & -- & $\circ$\textsuperscript{d} & -- & $\circ$ \\
Layer-wise SSL \cite{elkheir25_layerwise}
  & -- & -- & -- & -- & \checkmark \\
RADAR, comms.\ \cite{radar26,realworld_comms24}
  & -- & -- & -- & \checkmark & $\circ$ \\
WhisperKit (ASR) \cite{whisperkit25}
  & n/a & \checkmark & \checkmark & n/a & \checkmark \\
\midrule
\textbf{This work} & \checkmark & \checkmark & \checkmark & \checkmark & \checkmark \\
\bottomrule
\end{tabular}}
\\[2pt]
\parbox{\columnwidth}{\footnotesize
\textsuperscript{a}\,minDCF is primary; actDCF and $C_{\mathrm{llr}}$ are
reported. \textsuperscript{b}\,Codec variability is absorbed into domain shift,
not controlled. \textsuperscript{c}\,Browser on a desktop x86 CPU, not mobile
hardware. \textsuperscript{d}\,Compression by pruning, not quantization.}
\end{table}

\subsection{Efficient and edge detection}
A smaller body of work optimizes countermeasures for size and speed. Pascu et
al.\ \cite{pascu25_edge_plugin} truncate a 300M-parameter SSL model to its lower
layers (\(\sim\)\num{101}M) and run it inside a browser via ONNX Runtime;
layer-wise analyses \cite{elkheir25_layerwise} show that the earliest SSL layers
carry most of the spoofing-discriminative signal, motivating exactly such
truncation, and pruning has been applied to the same back-bones
\cite{pruning_ssl}. These systems shrink the model yet still evaluate it at an
oracle EER on a workstation --- and, in the deployed browser plugin, emit a hard
bona-fide/spoof verdict at a threshold that is never stated or measured. Their
axis is orthogonal to ours and the two compose: they hold the operating point
fixed and shrink the model, whereas we hold the model fixed and move the
operating point onto physical hardware. Compression that preserves oracle EER
tells us nothing about whether the threshold shipped with it still separates the
classes --- which is precisely the quantity we define and measure.
% REWRITTEN 2026-07-29 (D30). The previous version claimed "our audit *includes*
% a truncated SSL model rather than treating this work as external". That
% depended on roster slot #7, dropped at the roster freeze, so the claim is now
% false. This is the paper's weakest differentiation and reading_tiers.md flags
% pascu25 as "differentiation must survive review" -- do NOT restore the
% inclusion claim, and do not let this paragraph get thinner than it is.

\subsection{Robustness under channel degradation}
A parallel literature stresses the input signal. Recent challenges and studies
inject codec, VoIP, and real-world corruption to probe detector robustness
\cite{radar26,realworld_comms24}, and ASVspoof\,2021's deepfake track built
media-storage codec and compression variability directly into its evaluation set
\cite{yamagishi21_asvspoof2021}. This work and ours are complementary along a
clean axis: they perturb the \emph{input} while scoring at an oracle threshold,
whereas we hold the input fixed and perturb the \emph{operating point}. The two
are composable, and we adopt their corruption families inside our stress suite
while additionally freezing the threshold, so that a codec's effect surfaces as a
shift in deployed \(\frr/\far\) rather than as a change in oracle EER.

\subsection{On-device benchmarking and speech-model quantization}
Rigorous on-device measurement methodology exists, but for neighboring tasks.
MLPerf Mobile and Tiny \cite{reddi22_mlperf_mobile,mlperf_tiny} codify minimum
query counts, minimum run durations, tail-latency percentiles, and reporting
discipline for latency and energy, and recent systems
port this to on-device speech recognition and language models
\cite{whisperkit25}. Separately, compression of speech SSL models has been
studied for recognition: 8-bit quantization of wav2vec\,2.0 costs about
\SI{0.1}{\percent} WER \cite{pimentel23_ssl_compression}, with similar findings
for Whisper and mixed-precision speech foundation models
\cite{quantwhisper25,mixedprec25}. Two gaps remain. No open, reproducible
on-device benchmark exists for anti-spoofing, in contrast to the ASR and LM
tooling above. And every speech-compression study measures a threshold-free
quantity --- word error rate --- so the effect of quantization on a decision
\emph{threshold} has never been reported; because recognition has no threshold to
move, this absence is categorical rather than a matter of coverage. We port the
measurement discipline of the former and open the metric gap of the latter
with \qts{}.

% ============================ 3. PROTOCOL ==================================
\section{The \bench{} Protocol}
\label{sec:protocol}
% ~2 pp. Defend each choice in ONE sentence. Source: experiment_protocol.md.

\subsection{Score contract and threshold transfer}
% Label 1 = bona fide, score higher = bona fide, decision SPOOF if s<tau.
% tau calibrated on ASVspoof19 LA dev for EVERY model (D16 as collapsed by D30),
% then frozen (tau.py); analysis reads the frozen files, never re-derives (B3).
Each CM emits a scalar score \(s(x)\in\mathbb{R}\), higher meaning more
bona-fide-like; the decision is bona fide iff \(s(x)\ge\tau\). We fix
\(\taustar\) on the ASVspoof\,2019 LA development split --- the fine-tuning
source of every model in the roster --- at
three operating points and transfer it unchanged to every target, precision,
and stress condition:
\begin{itemize}
  \item \(\taustar_{\eer}\): the equal-error crossing on dev (interpolated);
  \item \(\taustar_{\frr1},\ \taustar_{\frr5}\): the strictest threshold whose
        dev bona-fide rejection is \(\le 1\%\) resp.\ \(\le 5\%\).
\end{itemize}
% NOTE the protocol's "FPR-x%" == this project's FRR (bona fide flagged spoof).

\subsection{Reported quantities}
For a transferred threshold \(\taustar\) on a target set we report
\(\hter(\taustar)=\tfrac12[\frr(\taustar)+\far(\taustar)]\) with the
\(\frr/\far\) decomposition mandatory, alongside the oracle \(\eer\) as
reference.

\subsection{Threshold transfer gap (\ttg{})}
% D19 candidate: reference term = min-HTER, NOT EER, so TTG is well-defined at
% all three operating points and is not a renaming of an EER-point identity.
\begin{equation}
\ttg(\taustar) \;=\; \hter(\taustar) \;-\; \min_{\tau}\hter(\tau),
\label{eq:ttg}
\end{equation}
the excess deployment error from transferring \(\taustar\) rather than an
oracle-optimal threshold on the target.% At the EER point the reference reduces
% to EER; we use min-HTER so TTG is defined at frr1/frr5 too. (see D19)

\subsection{Quantization-induced threshold shift (\qts{})}
% experiment_protocol.md sec.3.5. Needs BOTH the fp32 tau and the
% precision-matched tau -> recalibrate tau per precision on dev.
\begin{equation}
\qts \;=\; \hter\!\left(s_{q},\,\taustar_{\text{fp32}}\right)
        \;-\; \hter\!\left(s_{q},\,\taustar_{q}\right),
\label{eq:qts}
\end{equation}
for quantized scores \(s_q\), measured in-domain and on transfer targets, per
precision \(q\in\{\text{fp16},\text{int8}\}\).

\subsection{Sensitivity, stress suite, and device methodology}
% One para each: threshold-sensitivity curve + safe-band width; codec/noise/IR
% stress (applied to bona fide AND spoof); device protocol (thermal gate; MLPerf
% single-stream run rules adopted in D23: >=1,024 queries and >=60 s per cell,
% p50/p90/p95 reported; "warmup" is not an MLPerf concept, so do not claim it).
\placeholder{sensitivity / stress / device methodology}

% ============================ 4. PIPELINE & ROSTER ========================
\section{End-to-End Pipeline and Model Roster}
\label{sec:pipeline}
% ~2 pp. Roster table + selection rationale; conversion pipeline (runtimes,
% precisions); preprocessing-parity methodology; Android app (one figure).

\begin{table*}[!t]
\caption{Model roster (frozen, D30) and selection rationale. One representative
per int8 failure \emph{mechanism} rather than per architecture family. Params,
input, frontend, and predicted int8 outcome per
\texttt{reference\_architectures.md}.}
\label{tab:roster}
\centering
\begin{tabular}{llrllll}
\toprule
\# & Model & Params & Input & Frontend & Mechanism under test & Score conv.\ (author repo) \\
\midrule
1 & AASIST~\cite{jung22_aasist}           & 297K  & raw       & SincConv & graph attention & \verb|out[:,1]| \\
2 & AASIST-L~\cite{jung22_aasist}         & 85K   & raw       & SincConv & \emph{within-class control} & \verb|out[:,1]| \\
3 & RawNet2~\cite{tak21_rawnet2}          & 17.6M & raw       & SincConv & recurrence (GRU) & \verb|out[:,1]| \\
4 & Res-TSSDNet~\cite{hua21_res_tssdnet}  & 350K  & raw (6\,s)& conv     & plain convolution & \verb|out[:,0]-out[:,1]| \\
5 & XLSR+AASIST~\cite{tak22_xlsr_aasist}\,\textsuperscript{ref} & 315M & raw & SSL CNN & transformer & \verb|out[:,1]| \\
\bottomrule
\end{tabular}
\end{table*}
% NOTE per-repo score conventions differ (D9); Res-TSSDNet has bonafide=index 0.
% XLSR+AASIST is the UNRANKED server reference (D16), is in a different score
% space from zhou26, and enters reproduction UNGATED -- Tak et al. publish no
% 19LA-eval anchor. RawNet2's GRU is 16.1M of 17.6M params (92%), which is why
% an int8 row that falls back on it saves almost nothing -- state that in 6.1.
% AASIST/AASIST-L is the study's ONLY within-class replication; H6 is therefore
% stated at op level and declared hypothesis-generating, not predictive (D30).
% Excluded at freeze: RawGAT-ST, LFCC-LCNN, WavLM-trunc, WavLM+MHFA, pruned SSL.
% The three declared roster limitations (no handcrafted-feature model; one GAT
% lineage; ranked roster is pre-SSL-era) belong in Sec.~\ref{sec:discussion}.

\subsection{Conversion and quantization}
% PTQ config as methodology, not defaults (D17): S8S8, per-channel weights, QDQ,
% 500-clip calibration. Every int8 row ships its fp32-fallback op list (=H6).
% fp16 via graph conversion, not net.half() (D18). Preview conversion failures.
\placeholder{conversion pipeline; PTQ config; fp16 policy}

\subsection{Preprocessing-parity methodology}
% P3 PRIMARY GATE: zero decision flips at all three operating points between
% on-device feature extraction and the Python reference. Feature drift (1e-3)
% is diagnostic, not a gate.
\placeholder{parity methodology + Android app figure}

% ============================ 5. EXPERIMENTAL SETUP =======================
\section{Experimental Setup}
\label{sec:setup}
% ~1 p. Datasets/splits; phones table; seeds where retrained; CIs;
% pre-registration statement + repo link.

\subsection{Datasets}
% Four sets: ASVspoof19 LA (source; tau on dev), 21 DF, In-the-Wild, MLAAD.
% ASVspoof 5 is deferred (D30); D24 is dormant. Record ITW caveat: social-media
% provenance (tandem compression).
% 21 DF = eval subset on the D34 VERIFIED-AUDIO surface: 455,508 of 533,928
% trials. 78,420 are excluded because their only decode fails the file's own
% STREAMINFO MD5; the clean-only surface (305,525) is the sensitivity analysis.
% Published DF21 numbers (B04 22.38%, Tak et al. 2.85%) are NOT comparable: the
% organisers scored different audio for 43% of trials (df21_flac_integrity.md).
% OPEN [U]: make this a named data note?
\placeholder{datasets + splits table}

\begin{table}[!t]
\caption{Devices. \todo{fill from phones.md}}
\label{tab:phones}
\centering
\begin{tabular}{llll}
\toprule
Tier & Phone & SoC / RAM & OS build \\
\midrule
Flagship & \todo{} & \todo{} & \todo{} \\
Mid      & \todo{} & \todo{} & \todo{} \\
Low      & \todo{} & \todo{} & \todo{} \\
\bottomrule
\end{tabular}
\end{table}

\subsection{Hypotheses and statistical protocol}
% research_methodology.md Part C is BINDING. Exploratory vs confirmatory split;
% Holm-Bonferroni on confirmatory families; two cost models; 3-session device
% measures; decision-flip parity gate.
\begin{table}[!t]
\caption{Pre-registered hypotheses (audit\_plan.md). E = exploratory (effect
size + CI only); C = confirmatory (Holm--Bonferroni family).}
\label{tab:hypotheses}
\centering
\begin{tabular}{clc}
\toprule
H & Statement & Type \\
\midrule
H1 & Oracle EER anti-correlates with \ttg{} (target fixed) & E \\
H2 & \(\frr/\far\) asymmetry under transfer & C \\
H3 & \qts{}: quantization moves the operating point & C \\
H4 & \ttg{} vs parameter count (85\,K--315\,M) & E \\
H5 & Neural-codec (EnCodec) corrupts bona fide & C \\
H6 & int8 coverage is non-uniform across op classes & C \\
\bottomrule
\end{tabular}
\end{table}
% H1 CAUTION: test with the TARGET HELD FIXED -- zhou26's two points (LA->ITW
% TTG 28.3 vs LA->DF21 TTG 9.1, same model) show target prior dominates.
% H4 RESTATED by D30: "SSL scale buys separability, not calibration" needed
% multiple SSL points; the frozen roster has one and it is unranked. Consequence
% to own in Sec.~\ref{sec:discussion}: this study can no longer refute "just use
% a bigger model" with data -- answer it by argument and citation instead.
% H6 RESTATED at op level by D30. Declare the inferential limit in the text:
% with one model per op class this is hypothesis-generating, not a predictive
% rule; AASIST/AASIST-L is the only within-class replication we have.
% "Majority of models" in H2/H3/H5 = majority of the FOUR RANKED models;
% XLSR+AASIST is excluded from those counts (D16).
We pre-register H1--H6 \todo{repo commit hash} before any Stage-P5 run.

% ============================ 6. RESULTS ==================================
\section{Results}
\label{sec:results}
% ~4.5 pp. Subsections map to RQ1-3 + stress + synthesis (paper_outline sec.6).

\subsection{Deployability (RQ1)}
\label{sec:res-deploy}
% Conversion/parity matrix; op-level failure analysis (H6). "N/A -- does not
% fit" rows for the 315M SSL reference (inference graph 315.4M). This is where
% the H6 op inventory from the Tier-2 notes lands (ScatterND/topk in AASIST; GRU
% in RawNet2; SincConv constant-folds; TSSDNet clean control; XLS-R = transformer
% + grouped weight-normed Conv1d positional embedding + the AASIST back-end. It
% has NO relative-position-bias gather: that was WavLM's, and WavLM left the
% roster at D30. Verified in fairseq wav2vec2.py @a5402130, 2026-09-22).
\placeholder{conversion/parity matrix + op-level failure analysis}

\subsection{Operating behavior (RQ2)}
\label{sec:res-opbehav}
% TTG headline figure; FRR/FAR heatmap (H2); sensitivity curves; H1 scatter
% (oracle EER vs TTG); post-hoc calibration failure rows.
% INTEGRITY GATE (D20): any monotone transform (temperature scaling, z-norm)
% row reporting non-zero delta-EER is a bug/leak, not a finding (zhou26 Prop 1).
\placeholder{TTG figure, FRR/FAR heatmap, sensitivity, H1 scatter}

\subsection{Quantization (H3)}
\label{sec:res-quant}
% QTS table; decision flips; per-architecture sensitivity (SincConv vs conv vs
% attention). Watch int8 score-tie collapse coarsening the dev grid.
\placeholder{QTS table + decision flips}

\subsection{Cost (RQ3)}
\label{sec:res-cost}
% Pareto bubbles (HTER@tau* vs latency vs energy, size=params) per device tier;
% thermal throttling curves; scale-vs-calibration bars (H4).
\placeholder{Pareto bubbles, thermal curves}

\subsection{Channel stress}
\label{sec:res-stress}
% Codec curves incl. EnCodec bona-fide corruption (H5); noise/IR; MLAAD
% multilingual slice. Report WORST-CASE over conditions, not just means (D21).
\placeholder{codec/noise/IR curves}

\subsection{Synthesis: the deployability decision table}
\label{sec:res-synth}
% The practitioner artifact: which model to ship per device tier and use case,
% and what breaks.
\placeholder{deployability decision table}

% ============================ 7. DISCUSSION ===============================
\section{Discussion and Limitations}
\label{sec:discussion}
% ~1 p. PTQ-only (no QAT, out of scope); 3 phones (named case studies, not tier
% estimates; harness released); checkpoints as-is (removes retraining confound,
% inherits authors' training choices); recalibration guidance; ethics --
% fixed-threshold FRR failures disproportionately reject real speakers.
% Also concede: bona-fide-domain coverage is thin (mostly read speech +1
% conversational target); ITW confounds (parody spoofs, 8-bit provenance).
%
% ROSTER LIMITATIONS -- MANDATORY, from D30. Do not soften these:
%  (a) No handcrafted-feature model (LFCC-LCNN excluded), so on-device feature
%      parity is demonstrated only for raw-waveform preprocessing.
%  (b) One GAT lineage, so family-level generalization about graph-attention
%      models is not supported by this study's own data. AASIST/AASIST-L is the
%      only within-class replication.
%  (c) The ranked roster is pre-SSL-era (2021-22); the single SSL model is
%      unranked and expected to yield no device rows. NOTE this is structurally
%      the same criticism our own reading pass levelled at muller22, whose 12
%      architectures contained zero SSL models -- say so before a reviewer does.
%  (d) ASVspoof 5 deferred out of scope; the three shift axes (ITW corpus,
%      MLAAD language, 21DF+stress channel) are what carry the shift claims.
%  Defence, stated plainly: the contribution is the REGIME and the PROTOCOL,
%  which are roster-independent, plus a released harness that makes adding a
%  model cheap. Breadth is explicitly not the claim.
\placeholder{discussion + limitations}

% ============================ 8. CONCLUSION ===============================
\section{Conclusion}
\label{sec:conclusion}
% ~0.5 p. Restate the two findings; the harness as the durable contribution;
% the recommendation (report HTER at a transferred threshold alongside EER).
\placeholder{conclusion}

% ============================ ACK / BIB ===================================
\section*{Acknowledgment}
\placeholder{acknowledgments}

\bibliographystyle{IEEEtran}
\bibliography{refs}

% \begin{IEEEbiography}{...}\end{IEEEbiography}  % optional author bios

\end{document}
